\subsubsection{Schritt 3: Non-Maximum Suppression}

Schritt~2 liefert für jeden Pixel einen Gradientenbetrag, der die lokale
Kantenstärke beschreibt, wobei die resultierenden Kanten an vielen
Stellen breiter als ein Pixel sind. Non-Maximum Suppression soll diese
breite Gradientenantwort im Idealfall auf eine Breite von einem Pixel
reduzieren, indem entlang der Gradientenrichtung nur das lokale Maximum
erhalten bleibt. Canny begründet diese Forderung mit dem Kriterium
präziser Lokalisierung: \textit{,,The points marked as edge points by
the operator should be as close as possible to the center of the true
edge``} \cite[S.~680]{canny1986}. Formal entspricht der gesuchte
Kantenpunkt einer Nullstelle der zweiten Richtungsableitung des
geglätteten Bildes entlang der Gradientenrichtung $n$
\cite[S.~691]{canny1986}, die aus der geglätteten Gradientenrichtung
geschätzt wird, da sie nicht vorab bekannt ist.

In der diskreten Implementierung wird dies approximiert, indem die
Gradientenmagnitude jedes Pixels mit den beiden direkten Nachbarn
entlang der auf vier Klassen ($0^\circ$/$45^\circ$/$90^\circ$/$135^\circ$)
quantisierten Gradientenrichtung verglichen wird. Ein Pixel überlebt
die Suppression genau dann, wenn seine Magnitude größer oder gleich
beiden Nachbarn ist, andernfalls wird er auf null gesetzt. Die
Ein-Pixel-Breite wird dabei nur bei einem strikten lokalen Maximum
tatsächlich erreicht, denn bei einem Plateau, also mehreren
benachbarten Pixeln mit identischer Magnitude, erfüllen aufgrund des
$\geq$-Vergleichs alle beteiligten Pixel gleichzeitig das
Überlebenskriterium, wodurch die Kante an dieser Stelle breiter als ein
Pixel bleibt.